% page 548 of the facsimile (image p-547 of the scan)
% block 160.0 x 246.3 mm ; left margin 8.0 mm
\pgc{-12.700mm}{0.000mm}{
\l{\cel{3}540}
\l{}
\l{\sou{sofistiko}. La arto di la sofisto. DEFIRS.}
\l{}
\l{\sou{sofito}. (\sou{arkitekturo}). La suba parto di plafono, di arkitravo,}
\l{\cel{3}e c., ornita ye rozfenestri, e c. - DEFIRS.}
\l{}
\l{\sou{soko}. (\sou{agrokultivo}). Larja fero-peco triangula, tranchiva, mun-}
\l{\cel{3}tita an la plugilo-parto qua glitas sur sulo, e qua, kande}
\l{\cel{3}onu enterigas lu per preso, produktas la sulko, per fendar}
\l{\cel{3}la ter-surfaco. - EF.}
\l{}
\l{\sou{soklo}. (\sou{tekn.}) Suportilo, sub skultura busto, statueto, sub}
\l{\cel{3}vazo, ornivo arkitekturala. - DEFIRS.}
\l{}
\l{\sou{sokursar}. (\sou{tran.}) Helpar (ulu) eskapar danjero qua tre minacas}
\l{\cel{3}lu. - EFIS.}
\l{}
\l{\textquotedbl{}\sou{sol}\textquotedbl{}. - (\sou{muziko}). Noto kinesma di la \textquotedbl{}do\textquotedbl{}-gamo.}
\l{}
\l{\sou{sola}. Qua esas ne-akompanata da sua simili ; qua ne esas kun}
\l{\cel{3}sua simili. - DEFIRS.}
\l{}
\l{\sou{solano}. (\sou{bot.)}\cel{2}Genero de planti ek la familio \textquotedbl{}solanei\textquotedbl{}. -}
\l{\cel{3}L.\cel{1}\sou{morchella esculenta}. IL.}
\l{}
\l{\sou{soldar}. (\sou{trans.})\cel{2}Kunjuntar (speci de metalo) per kompozajo}
\l{\cel{3}metala fuzebla (per shalmo). - (Cf.\cel{1}\sou{weldar}). - EFIS.}
\l{}
\l{\sou{soldato}. (\sou{milit.})\cel{2}Armeano sen-grada.- DEFIRS.}
\l{}
\l{\sou{soleo}. (\sou{zool.)}\cel{3}Genero de fishi, teleostea, familio ek la}
\l{\cel{3}pleuronekti, qui vivas en omna mari. - L.\cel{1}\sou{pleuronectes}}
\l{\cel{3}\sou{solea}. - (\sou{anat.}) Un ek la muskuli di la gambo. - EFI}
\l{}
\l{\sou{solecismo}. Kulpo pri la reguli di la sintaxo di linguo. DEFIRS}
\l{}
\l{\sou{solena}. I. Celebrata singla-yare per ceremonio publika. -}
\l{\cel{3}II. Akompanata da ceremonio, da agi publika, qui igas lu}
\l{\cel{3}partikulare importoza ed impozanta. - III. (\sou{metaf}.)Grava}
\l{\cel{3}til igar la kozo importoza ed impozanta. - DEFIS.}
\l{}
\l{\sou{solenoido}. (\sou{elektro}). Ensemblo de mikra ringi inter-egala,}
\l{\cel{3}inter-paralela, tre inter-proxima, egal-distanta, perpen-}
\l{\cel{3}dikla ye un rekto qua unionas lia centri, ed en qua fluas}
\l{\cel{3}elektro sama-sinse. - DEFIRS.}
\l{}
\l{\sou{solfejar}. (\sou{trans.})\cel{2}(\sou{muziko}).Lektar (frazo muzikala) intonante}
\l{\cel{3}e nomante singla noto, segun la movo-rapideso e la mezuro.}
\l{\cel{3}- DEFR.}
\l{}
\l{\sou{solio.}\cel{1}Peco de ligno o de quadro sur qua rezeskas tota-larjese}
\l{\cel{3}la infra parto di pordo, quan onu transiras por enirar, e}
\l{\cel{3}qua, lokizita super la nivelo di la sulo extera, posibligas}
\l{\cel{3}- per impedar la eniro subporda di aquo, tero, e c. - la}
\l{\cel{3}plu facila funciono di la pordo ed igas la klozo plu tota.}
\l{\cel{3}- FI}
}
